

\chapter{Wissenschaftliche Fragestellung}\label{cha:researchQuestion}

Diese Arbeit untersucht, wie gut sich MOMENT-1-small für eine cloudbasierte \ac{RUL}-Prognose mit fehlerhaften Sensordaten eignet. Als Vergleich dienen ein \ac{LSTM} und ein \ac{TCN}, die speziell für diese Aufgabe trainiert werden. Bewertet werden die Prognosequalität und der Aufwand für den Betrieb in der Cloud. Daraus ergibt sich folgende Forschungsfrage:

\begin{quote}
	Wie unterscheiden sich MOMENT-1-small, ein \ac{LSTM} und ein \ac{TCN} bei der \ac{RUL}-Prognose mit fehlerhaften \ac{CMAPSS}-Sensordaten hinsichtlich Prognosequalität, Robustheit, Inferenzlatenz, Ressourcenbedarf und Cloud-Kosten?
\end{quote}

\noindent Zur Beantwortung der Forschungsfrage werden folgende Teilfragen untersucht:

\begin{enumerate}
	\item Wie verändern sich die Prognosefehler der Modelle bei unterschiedlichen Arten und Stärken von Sensorfehlern?
	\item Welche Unterschiede zeigen sich bei Inferenzlatenz, Durchsatz und Ressourcenbedarf unter vergleichbaren Betriebsbedingungen?
	\item In welchem Verhältnis stehen die Kosten des Cloud-Betriebs zur erreichten Prognosequalität und Robustheit?
\end{enumerate}

Für die Hauptuntersuchung ist \ac{CMAPSS} FD001 vorgesehen. FD003 soll den Vergleich ergänzen. Das Training für die \ac{RUL}-Prognose erfolgt ohne zusätzlich erzeugte Sensorfehler. Erst die Testdaten werden verändert, um eine Verschlechterung der Datenqualität während des Betriebs nachzubilden. Die Modelle werden nicht nachträglich an diese Testfehler angepasst.
